\subsection{THE AXIOMS}

An equalitarian theory is a theory \(\mathcal{C}\) which has a relational sign of weight 2, written = (read ``equals''), and in which the schemes S1 through S5 (§§3 and 4), together with the schemes S6 and S7 below, provide implicit axioms. If \(T\) and \(U\) are terms in \(\mathcal{C}\) , the assembly = \(TU\) is a relation in \(\mathcal{C}\) (called the relation of equality) by virtue of CF4; in practice it is denoted by \(T = U\) or \((T) = (U)\) .

S6. Let x be a letter, let T and U be terms in C, and let \(R\{x\}\) be a relation in C; then the relation \((T = U) \Longrightarrow (R\{T\} \Longleftrightarrow R\{U\})\) is an axiom.

S7. If R and S are relations in C and if x is a letter, then the relation \(((\forall x)(R \leftrightarrow S)) \Longrightarrow (\tau_{x}(R) = \tau_{x}(S))\) is an axiom.

To show that the rule S6 is a scheme, let A be an axiom of G, obtained by applying S6; then there is a relation R in G, terms T and U in G, and a letter x, such that A is \((T=U)\Longrightarrow((T|x)R\Longleftrightarrow(U|x)R)\) . We shall show that if y is a letter and V a term in C, the relation \((V|y)A\) can be obtained by applying S6. By means of CS1 (§1, no. 2) we may assume that x is distinct from y and does not appear in V. Let \(T'\) , \(U'\) , \(R'\) denote the assemblies \((V|y)T\) , \((V|y)U\) , \((V|y)R\) respectively. By CS2 and CS5 (§1, no. 2), \((V|y)A\) is identical with

\[
(T ^ {\prime} = U ^ {\prime}) \Longrightarrow ((T ^ {\prime} | x) R ^ {\prime} \Longleftrightarrow (U ^ {\prime} | x) R ^ {\prime})
\]

and the proof is complete. The verification that S7 is a scheme is similar.

Intuitively, the scheme S6 means that if two objects are equal, they have the same properties. Scheme S7 is more remote from everyday intuition; it means that if two properties R and S of an object x are equivalent, then the distinguished objects \(\tau_{\mathbf{x}}(\mathbf{R})\) and \(\tau_{\mathbf{x}}(\mathbf{S})\) (chosen respectively from the objects which satisfy R, and those which satisfy S, if such objects exist) are equal. The reader will note that the presence in S7 of the quantifier \(\forall x\) is essential (cf.~Exercise 7).

The negation of the relation = TU is denoted by \(T \neq U\) or \((T) \neq (U)\) (where the sign \(\neq\) is read ``is different from'').

From S6 we deduce the following criterion :

C43. Let x be a letter, let T and U be terms in G, and let \(R\{x\}\) be a relation in G. Then the relations

\[
(\boldsymbol {T} = \boldsymbol {U} \text {   and   } \boldsymbol {R} \{\boldsymbol {T} \}), \quad (\boldsymbol {T} = \boldsymbol {U} \text {   and   } \boldsymbol {R} \{\boldsymbol {U} \})
\]

are equivalent.

For if we adjoin the hypotheses \(T = U\) and \(R\{T\}\) , then \(R\{U\}\) is true by S6; therefore \((T = U \text{ and } R\{U\})\) is true.

When a relation of the form \(T = U\) has been proved in a theory \(\mathfrak{C}\) , it is often said (by abuse of language) that \(T\) and \(U\) ``are the same'' or are ``identical''. Likewise, when \(T \neq U\) is true in \(\mathfrak{C}\) , we say that \(T\) and \(U\) ``are distinct'' in place of saying that \(T\) is different from \(U\) .

